% ═══════════════════════════════════════════════════════════════════
%  PART 3: Diagnostics, Python Usage & Interpretation
% ═══════════════════════════════════════════════════════════════════

\section{Pre-Estimation Diagnostic Tests}

Before applying KRLS/QQKRLS, researchers must \textbf{justify} the use of nonparametric methods. The \texttt{qqkrls} library includes a complete diagnostic battery.

\subsection{Ramsey RESET Test}

\begin{definitionbox}[Ramsey RESET (Regression Equation Specification Error Test)]
Tests whether the linear model is correctly specified by including powers of fitted values:
\begin{equation}
y = X\beta + \gamma_2 \hat{y}^2 + \gamma_3 \hat{y}^3 + \varepsilon
\end{equation}
\begin{itemize}[leftmargin=1.5em]
\item $H_0$: Linear model is correctly specified ($\gamma_2 = \gamma_3 = 0$)
\item $H_1$: Nonlinear terms are significant $\to$ \textbf{use KRLS}
\end{itemize}
\end{definitionbox}

\subsection{Breusch--Pagan Test}

\begin{definitionbox}[Breusch--Pagan Test for Heteroskedasticity]
Tests whether the variance of residuals depends on the predictors:
\begin{equation}
\hat{\varepsilon}^2 = X\alpha + u
\end{equation}
\begin{itemize}[leftmargin=1.5em]
\item $H_0$: Homoskedasticity ($\Var[\varepsilon \mid X] = \sigma^2$)
\item $H_1$: Heteroskedasticity $\to$ \textbf{effects differ across distribution} $\to$ use quantile-based analysis
\end{itemize}
\end{definitionbox}

\subsection{BDS Test}

\begin{definitionbox}[BDS Test (Brock, Dechert, Scheinkman)]
A nonparametric test for serial independence and nonlinear structure in the residuals:
\begin{itemize}[leftmargin=1.5em]
\item $H_0$: Residuals are i.i.d.
\item $H_1$: Nonlinear dependence exists $\to$ \textbf{justifies KRLS/QQKRLS}
\end{itemize}
The test uses the correlation integral at embedding dimension $m$:
\begin{equation}
C_m(\epsilon) = \lim_{N \to \infty} \frac{2}{N(N-1)} \sum_{i<j} \mathbf{1}\!\left[\|x_i^m - x_j^m\| < \epsilon\right]
\end{equation}
Under the null, $C_m(\epsilon) = [C_1(\epsilon)]^m$. The BDS statistic tests this equality.
\end{definitionbox}

\subsection{Jarque--Bera Normality Test}

\begin{definitionbox}[Jarque--Bera Test]
Tests whether the data follows a normal distribution:
\begin{equation}
JB = \frac{N}{6}\left(S^2 + \frac{(K-3)^2}{4}\right) \sim \chi^2(2)
\end{equation}
where $S$ = skewness, $K$ = kurtosis. Rejection suggests \textbf{non-Gaussian distributions}, motivating quantile-based methods.
\end{definitionbox}

\subsection{Andrews Parameter Stability Test}

\begin{definitionbox}[Andrews Sup-F Test]
Tests for structural breaks in the relationship by computing $F$-statistics at all possible break points:
\begin{itemize}[leftmargin=1.5em]
\item \textbf{Max-F}: Supreme of $F$-statistics (detects single break)
\item \textbf{Exp-F}: Exponential average (robust to multiple breaks)
\item \textbf{Ave-F}: Simple average (overall instability)
\end{itemize}
Parameter instability justifies quantile-based methods that allow coefficients to vary.
\end{definitionbox}


\section{Using the \texttt{qqkrls} Python Library}

\subsection{Installation}

\begin{lstlisting}[language=bash,style=pythonstyle]
pip install qqkrls
\end{lstlisting}

\subsection{Complete Workflow}

\begin{lstlisting}
import numpy as np
import pandas as pd
from qqkrls import (
    krls, qqkrls, predict_krls,
    plot_qqkrls_heatmap, plot_qqkrls_3d,
    plot_qqkrls_contour, plot_qqkrls_pvalue,
    plot_krls_derivatives, plot_krls_fit,
    descriptive_statistics, linearity_test_battery,
    bds_test, jarque_bera, print_diagnostics,
)

# Step 1: Load data
y = df['EF'].values
X = df[['Energy', 'CO2', 'GDP']].values

# Step 2: Pre-estimation diagnostics
diag = linearity_test_battery(y, X, col_names=['Energy','CO2','GDP'])
print_diagnostics(diag)

# Step 3: Fit KRLS
fit = krls(X, y, derivative=True, vcov=True, verbose=1)

# Step 4: QQKRLS estimation
taus = np.arange(0.05, 1.0, 0.05)
result = qqkrls(y=y, x=df['Energy'].values,
                y_quantiles=taus, x_quantiles=taus,
                n_boot=200, verbose=True)

# Step 5: Visualize
plot_qqkrls_heatmap(result, colorscale='paper',
                    show_stars=True, save_path='heatmap.png')
plot_qqkrls_3d(result, save_path='surface.png')
\end{lstlisting}

\subsection{Key API Functions}

\begin{center}
\small
\begin{tabular}{>{\ttfamily\color{deepblue}}l p{8cm}}
\toprule
\textbf{Function} & \textbf{Description} \\
\midrule
krls(X, y) & Fit KRLS model with LOO cross-validation \\
qqkrls(y, x) & Run QQKRLS with bootstrap inference \\
predict\_krls(fit, X\_new) & Out-of-sample prediction \\
\midrule
plot\_qqkrls\_heatmap() & Paper-style coefficient heatmap \\
plot\_qqkrls\_3d() & MATLAB-style 3D surface plot \\
plot\_qqkrls\_contour() & Filled contour plot \\
plot\_qqkrls\_pvalue() & P-value significance map \\
plot\_qqkrls\_panel() & Multi-variable panel \\
plot\_krls\_derivatives() & Marginal effect scatter + density \\
plot\_krls\_fit() & Fitted vs actual with $R^2$ \\
\midrule
descriptive\_statistics() & LaTeX descriptive statistics table \\
qqkrls\_coefficient\_table() & LaTeX coefficient matrix \\
krls\_summary\_table() & LaTeX KRLS summary \\
export\_results\_csv() & CSV export of all results \\
\midrule
linearity\_test\_battery() & RESET + BP + JB + BDS \\
bds\_test() & BDS independence test \\
parameter\_stability\_test() & Andrews sup-F, exp-F, ave-F \\
jarque\_bera() & Jarque--Bera normality test \\
\bottomrule
\end{tabular}
\end{center}


\section{Interpreting QQKRLS Results}

\subsection{Reading the Heatmap}

The QQKRLS heatmap is the primary output. Each cell $(\theta, \tau)$ represents one nonparametric marginal effect:

\begin{warningbox}[Common Misinterpretations to Avoid]
\begin{itemize}[leftmargin=1.5em]
\item The axes are \textbf{quantiles}, not raw values of the variables.
\item $\theta = 0.05$ means ``when $X$ is at its 5th percentile,'' not ``when $X = 0.05$.''
\item A positive coefficient means: at that distributional location, increasing $X$ \emph{increases} $Y$.
\end{itemize}
\end{warningbox}

\subsection{Pattern Recognition}

\begin{center}
\begin{tabular}{>{\bfseries\color{deepblue}}l p{9cm}}
\toprule
\textbf{Pattern} & \textbf{Economic Interpretation} \\
\midrule
Uniform positive & $X$ always increases $Y$ (homogeneous positive effect) \\
Gradient left-to-right & Effect strengthens as $X$ increases (nonlinear amplification) \\
Gradient bottom-to-top & Effect stronger when $Y$ is high (asymmetric vulnerability) \\
Diagonal pattern & Effect strongest at matching quantiles (distributional symmetry) \\
Quadrant split & Effect changes sign across the distribution (regime switching) \\
Significance clusters & Only certain quantile pairs show significant effects \\
\bottomrule
\end{tabular}
\end{center}

\subsection{Writing Up Results for a Paper}

\begin{examplebox}[Template for Results Section]
``Table X presents the QQKRLS estimates of the marginal effect of [Variable] on [Dependent Variable] across 19 quantile pairs ($\theta, \tau \in \{0.05, 0.10, \ldots, 0.95\}$). The pre-estimation diagnostics (Ramsey RESET: $F = \ldots$, $p = \ldots$; BDS: $z = \ldots$, $p < 0.01$) reject linearity and i.i.d.\ residuals, justifying the nonparametric QQKRLS approach.

The heatmap in Figure Y reveals substantial heterogeneity across the distribution. At lower quantiles of [Variable] ($\theta < 0.30$), the effect on [DV] is [negative/insignificant], while at upper quantiles ($\theta > 0.70$), the effect becomes [significantly positive]. This asymmetry suggests that [economic interpretation].

Of the 361 quantile-pair cells, [N] ([\%]\%) are significant at the 5\% level, concentrated in the [describe region] of the heatmap. These findings are robust to the bootstrap inference with $B = 200$ replications.''
\end{examplebox}

\subsection{Methodological Justification Workflow}

The recommended workflow for any QQKRLS paper follows five steps:

\begin{center}
\begin{tikzpicture}[
  node distance=1.2cm,
  box/.style={rectangle, rounded corners=6pt, draw=primaryblue, fill=primaryblue!8, text width=9cm, minimum height=1.2cm, align=center, font=\small},
  arrow/.style={-{Stealth[length=6pt]}, thick, primaryblue}
]
\node[box] (s1) {\textbf{Step 1:} Run diagnostics (RESET, BP, BDS, JB) $\to$ Reject linearity};
\node[box, below=of s1] (s2) {\textbf{Step 2:} Fit KRLS $\to$ Confirm high $R^2$, check residuals};
\node[box, below=of s2] (s3) {\textbf{Step 3:} Examine ME heterogeneity (Q25 $\neq$ Q75) $\to$ Justify quantile approach};
\node[box, below=of s3] (s4) {\textbf{Step 4:} Run QQKRLS with bootstrap $\to$ Coefficient matrix $\beta(\theta, \tau)$};
\node[box, below=of s4] (s5) {\textbf{Step 5:} Visualize (heatmap, 3D, contour) $\to$ Interpret patterns};
\draw[arrow] (s1) -- (s2);
\draw[arrow] (s2) -- (s3);
\draw[arrow] (s3) -- (s4);
\draw[arrow] (s4) -- (s5);
\end{tikzpicture}
\end{center}
